% SPDX-License-Identifier: Apache-2.0
% covers: Uart
\datasheet{Uart}{Vreteno's serial port: an AXI-Lite peripheral with a line each way and a receive buffer of eight.}

\dsfacts{\code{cpu/vreteno/src/uart.rs}}{\code{vreteno32::uart::Uart}}%
{\code{//docs:vreteno}, \code{//docs:noc}}

\subsection*{Function}

\code{Uart} sends bytes on \code{tx} and receives them on \code{rx},
one start bit, eight data bits least significant first and one stop
bit, each \code{DIV} cycles long. The lines rest high. A received byte
lands in a buffer of eight, and \code{irq} is high while the buffer
holds any.

It is an AXI-Lite peripheral. It sits behind a \code{LiteBridge1},
which hands it one transaction at a time with no identifier and no
burst. Reads come on \code{ar} and are answered on \code{r}; writes
come on \code{aw} and \code{w} and are answered on \code{b}. The input
\code{rst} stops both sides and empties the buffer.

\subsection*{Parameters}

\begin{center}\footnotesize
\begin{tabular}{@{}lp{0.7\textwidth}@{}}
\toprule
Parameter & Meaning \\
\midrule
\code{DIV} & Cycles per bit on both lines. The tables use 868, which
\code{uart.rs} states is 115200 baud at 100~MHz, as the board runs.
The checked runs use 4, and \code{board\_netlist sim} uses 16. \\
\bottomrule
\end{tabular}
\end{center}

\subsection*{Register map}

Offsets from the port's base. On the board the base is \code{0x3000},
which is also \code{isa::UART\_BASE}. The port selects the word by
bits 3 and 2 of the address.

\begin{center}\footnotesize
\begin{tabular}{@{}lp{0.38\textwidth}p{0.38\textwidth}@{}}
\toprule
Offset & Read & Write \\
\midrule
\code{0x0} & The last byte accepted for sending, zero-extended. &
Bits 7 to 0 start a frame if the port is idle; else the byte is
dropped. \\
\code{0x4} & Status: bit 0 busy sending, bit 1 a byte received and not
yet read, bit 2 the buffer full. Other bits zero. & Ignored. \\
\code{0x8} & The oldest received byte, zero-extended. The read removes
it from the buffer. & Ignored. \\
\code{0xc} & Zero. & Ignored. \\
\bottomrule
\end{tabular}
\end{center}

\dstables{Uart}

\subsection*{Behaviour}

\begin{itemize}
\item A read is taken when \code{r} has room and a request is on
\code{ar}, and it is answered with \code{OKAY} in the same cycle.
\item A write is taken when \code{b} has room and both \code{aw} and
\code{w} hold a request. It is answered with \code{OKAY} in the same
cycle, whichever word it names.
\item A byte written to word 0 while the port is idle loads the ten
bits of its frame into \code{shift} and ten into \code{bits}. It also
goes into \code{last}, and \code{sent} counts it.
\item While \code{bits} is not zero the port is busy and \code{tx} is
bit 0 of \code{shift}. Every \code{DIV} cycles the frame shifts down
and \code{bits} counts down. When idle, \code{tx} is high.
\item The \code{rx} line passes through one register, \code{line},
before the logic.
\item When the port is not receiving and \code{line} is low, a frame
begins. The port samples \code{line} in cycle \code{DIV/2} of each of
the ten bits.
\item A high where the start bit should be is a false start, and the
frame is dropped. The eight data bits shift into \code{rx\_shift}.
\item A high stop bit lands the byte. If the buffer holds eight, the
byte is dropped and \code{dropped} counts it; else it goes in at the
tail and \code{received} counts it. A low stop bit lands nothing.
\item A read of word 2 with a byte in the buffer removes the oldest. A
push and a removal in the same cycle leave the count unchanged.
\item \code{irq} is high while the buffer's count is not zero.
\item With \code{rst} high, \code{bits}, \code{tick},
\code{rx\_bits}, \code{rx\_tick}, \code{head} and \code{count} go to
zero.
\end{itemize}

\subsection*{Verification}

\begin{itemize}
\item \code{//cpu/vreteno:lockstep} puts the terminal of
\code{term.rs} on the two lines. At the halt it checks \code{sent} and
\code{last} against the bytes the model was given, and that
\code{dropped} is zero. In \code{demo\_program} the port must say
\code{OK} and a newline, receive three bytes, and echo them.
\item \code{//cpu/vreteno:hello\_test},
\code{//cpu/vreteno:hello\_cc\_test} and \code{//cpu/vreteno:board\_test}
check what compiled programs print on \code{tx}.
\item The build simulates \code{Uart<4>}'s netlist, the entity
\code{uart4}, against the trace of \code{//cpu/vreteno:vreteno} under
nvc and Verilator: \code{//docs:vreteno\_sim\_uart4\_uart4\_tb\_test}
and \code{//docs:vreteno\_vsim\_uart4\_test}. The run also writes
\code{Uart<868>}'s netlist, which no test simulates.
\item \code{//docs:vreteno} reports that the board, programmed with the
design from before the DDR3 memory was on the bus, said \code{OK} at
115200 baud and echoed the three bytes a watcher typed.
\end{itemize}

\subsection*{Used by}

\code{Board}, as \code{uart} behind the bridge \code{puart}. The
demonstration \code{//cpu/vreteno:vreteno}, the lockstep test and
\code{run.rs}. The system in \code{soc}, where it is the serial port at
corner 1,1 of the lattice.

\subsection*{Limits}

A byte written while a frame is going out is dropped, and a byte
received with the buffer full is dropped and counted. A program polls
the status and keeps up. \code{sent}, \code{received} and
\code{dropped} are eight-bit counters that are not on the bus. The port
checks no address beyond bits 3 and 2, since the bridge sends it only
its own range. The frame is fixed at ten bits, with no parity bit, and
the port has no flow-control lines.
